One way to mitigate poor conditioning is to change the geometry of the update.
Instead of \eqref{eq:gd-update}, consider
\begin{equation}
  x^{(k+1)} = x^{(k)} - t^{(k)} A \nabla f(x^{(k)}),
  \label{eq:preconditioned-gd}
\end{equation}
where $A \succ 0$ is a chosen matrix (a \emph{preconditioner}).
Equivalently, \eqref{eq:preconditioned-gd} corresponds to a linear change of variables that aims to make level sets of $f$ ``more spherical'' so that one step size works well in all directions.
Choosing a good $A$ typically requires problem-specific structure; see \citet[\S9.4]{boyd_vandenberghe_2004_convex}.

If $A$ is positive definite, the direction $-A\nabla f(x)$ is the steepest-descent direction when distance is measured in the norm $\norm{u}_{A^{-1}}=\sqrt{u^\top A^{-1}u}$.
Good choices of $A$ make this geometry align with the curvature of $f$; heuristically, $A \approx \Hess f(x)^{-1}$ reduces the effective condition number.
For the ridge logistic objective \eqref{eq:logistic-ridge-objective}, this picture is quite concrete.
The weighted Gram matrix $X^\top W(\theta)X$ records how sharply the loss curves in different directions, while the ridge term $mI$ adds the same amount of curvature everywhere.
A useful preconditioner therefore tries to match this combined geometry rather than treating every coordinate direction identically.

In practice, explicit matrix preconditioners are often approximated or restricted to be cheap.
For example, many ``adaptive'' optimizers in machine learning can be interpreted as choosing a diagonal (or block-diagonal) preconditioner based on running second-moment statistics of the stochastic gradients; see the EMA discussion in \cref{subsubsec:sgd-practical-guidance}.

